% =====================================================================
%  CHƯƠNG 6. KẾT LUẬN VÀ HƯỚNG PHÁT TRIỂN
% =====================================================================
\chapter{KẾT LUẬN VÀ HƯỚNG PHÁT TRIỂN}
\label{chap:6}

% ---------------------------------------------------------------------
\section{Tóm tắt kết quả đạt được}
\label{sec:6.1}

\canviet{Tổng kết theo từng chương: cơ sở lý thuyết và chứng minh
($G\HT = 0$, $\dmin = 3$, mã hoàn hảo, mã (8,4)); ví dụ tính tay; phần mềm mô
phỏng; kết quả kiểm thử và mô phỏng BSC. Chỉ nêu kết quả đã có số liệu ở
Chương 5.}

% ---------------------------------------------------------------------
\section{Hạn chế và hướng phát triển}
\label{sec:6.2}

\canviet{Hạn chế: chỉ sửa được 1 lỗi mỗi khối; phần mềm (7,4) không báo được
lỗi kép; tỉ lệ mã $4/7$. Hướng phát triển: cài mã (8,4) vào phần mềm; mã
Hamming tổng quát $(2^m-1,\,2^m-1-m)$; xen kẽ bit (interleaving) chống lỗi
chùm; so sánh với các mã khác.}
